% Source: 01 Properties_of_Matter 26(1).docx; formatting in layout.tex.
\HOne{Classification of Substances}

\HTwo{Mixtures, Pure Substances, Compounds and Elements}

Matter can be classified into different categories: mixtures or pure substances


A pure substance like distilled water is composed of only a single type of particles. It cannot be separated into different substances by physical methods.


However, pure substances can be divided into two different categories depending on whether or not they can be broken down into simpler substances by chemical reactions.


The helium in a party balloon is a good example of an \textbf{element}. It cannot be broken down into simpler substances by chemical reactions.


Water on the other hand is a good example of a \textbf{compound}. Even though it has been considered an element for hundreds of years, it can be broken down into simpler substances by chemical reactions.


\begin{center}% Replaces images/image43.jpeg (original retained)
\resizebox{\linewidth}{!}{\PureSubstanceClassificationFigure}\end{center}


\HTwo{Elements and Compounds}

The concept of\textbf{ elements} is ancient. It developed in different civilizations in an attempt to explain the huge variety of substances which exist. It also provided an understanding of the nature of change which for example occurs when a piece of wood rots, or is burnt to produce charcoal or ash. Most widely known are the four elements "earth, air, fire and water" that were popularized by Greek philosophers between 500 and 400 BC.


The beginnings of scientific experimentation in the 18th century led to our modern concept of the chemical element. It defines an element as a pure chemical substance which cannot be separated into simpler substances by chemical reactions.


Today around 118 chemical elements are known. Ninety-four of them have been found in nature. Around twenty more have been made artificially.


The elements are the building blocks from which all substances are made.


Chemical symbols are used as abbreviations for the elements. They consist of one or two letters, with the first letter, but not the second, capitalized.


They include  \ce{H} (\rule{35mm}{0.4pt}),  \ce{O} (\rule{35mm}{0.4pt}),


\ce{C} (\rule{35mm}{0.4pt}), and \ce{Ag} (\rule{35mm}{0.4pt}).


Most elements can interact with other elements to form \textbf{compounds}. When hydrogen reacts with oxygen, for example, water is formed. On the other hand, water can be broken down into oxygen and hydrogen by passing an electrical current through it.


The properties of water are completely different from the properties of hydrogen and oxygen. Hydrogen, oxygen and water are each unique substances with unique properties.


Compounds are pure chemical substances consisting of two or more different chemical elements. They can be separated into simpler substances (compounds or elements) by chemical reactions.


More than 20 million different compounds are known today. This means that more than 99.99 \% of all pure substances are compounds.


Chemical formulas are used to represent compounds such as H2O for \rule{35mm}{0.4pt},


\ce{NaCl} for \rule{35mm}{0.4pt}, \ce{H2SO4} for \rule{35mm}{0.4pt},


\ce{NH3} for \rule{35mm}{0.4pt}, and \ce{Ag2O} for \rule{35mm}{0.4pt}.


\HTwo{Summary}

\HFour{Physical Processes}
\begin{center}% Replaces images/physical_summary.png (original retained)
\resizebox{0.92\linewidth}{!}{\PhysicalProcessSummaryFigure}\end{center}


\HFour{Chemical Reactions}
\begin{center}% Replaces images/chemical_summary.png (original retained)
\resizebox{0.92\linewidth}{!}{\ChemicalReactionSummaryFigure}\end{center}

